\chapter{Growing Pains of Industrialization}

Throughout 1927 opinion in the party moved steadily in favour of rapid industrialization and the five-year plan, though hostile or sceptical reactions still made themselves heard, and the full consequences of these ambitious designs were not faced or understood. The expulsion of the opposition at the party congress of December 1927 cleared the ground by making it possible to silence criticism, and to adopt without undue embarrassment policies which the opposition had championed in the past. The grain collections crisis which followed the congress speeded the process. The first condition of industrialization was that the peasant should supply the necessary food for towns and factories at prices which did not put an intolerable strain on wage levels, and without the diversion of more than a minimum of the resources of industry to the manufacture of consumer goods for the peasant market. These problems had bedevilled the grain collections after the harvest of 1927, and seemed at first insoluble. The successful operation of the first months of 1928 was interpreted as showing that, with a sufficient application of force, the coercion of the peasant was practicable as well as indispensable. The peasant, the recalcitrant factor in a planned socialist economy, had been tamed. During 1928 inhibitions were gradually overcome, and industrialization driven forward relentlessly. The way was open. To force the pace required only an iron will to surmount obstacles, if necessary, by the same methods of coercion. Heroic determination and callous brutality were both manifested in the process. 
% 注：去除了 "[for towns" 和 "fable strain" 等侧边乱码。

The tensions set up by forced industrialization extended far beyond the world of the peasant. The revolution had installed new men in the seats of power. But it had not had time to rear and train a new generation of officials, experts, scientists, industrial managers, engineers and technicians of all kinds, whose services were indispensable to any regime; and these services were still performed mainly by the same men who had performed them for the last Tsar and for the Provisional Government. The group of officials and experts who manned the People's Commissariats and other Soviet institutions also included a considerable number of ex-Mensheviks and ex-SRs; the former predominated in Gosplan and Narkomfin, the latter in Narkomzem. Most of these non-party servants of the regime, who had reconciled themselves to the principles of NEP, profoundly disliked and mistrusted the new policies, advised against them, and showed no alacrity---sometimes, perhaps, passive resistance---in carrying them out. From this it was a short step to suspicions, widely entertained in party circles, of an active conspiracy to sabotage them. The wholesale dismissal of non-party officials and experts from positions of influence in Narkomfin and Narkomzem, the two commissariats which put up the most stubborn resistance to forced industrialization, began in the spring of 1928. The most sensational episode occurred in March, when 55 engineers and administrators employed in the coal-mines of the Don basin were arrested on a charge of sabotage, allegedly organized from abroad. After a mass public trial, in which many of the defendants made confessions, eleven death sentences were pronounced, and five were carried out. Others received long prison sentences. Three German engineers originally charged with complicity were acquitted. The trial set a pattern for future demonstrative denunciations and show trials. But for the moment the authorities recoiled somewhat from the suspicion and hostility generated against ``specialists'' of bourgeois background who were essential to the maintenance and expansion of industry, and issued several reassuring pronouncements. The training of workers as qualified engineers proceeded slowly; and the employment of foreign engineers, at first mainly German, later American, on major construction works was a feature of the period. 
% 注：去除了 "[far beyond", "1", "1/", "/stubborn", "f public trial", "/The trial", "[and show", ".“specialists”" 等诸多乱码。修复了断行错位 "of sabotage, allegedly organized from abroad. After a mass public trial..."。

It was not only the management side of industry which was subjected to the growing pressure. If the first, though still unstated, condition of industrialization was that the peasant should deliver his grain to the towns for a modest return, the second, and openly avowed, condition was that the productivity of the worker should increase faster than his wages, so that industrial expansion could be in part financed out of the profits of industry itself---an alternative to the unbridled exploitation of the peasant. This had been the chief aim of the campaigns of 1926 and 1927 for the ``regime of economy'' and ``the rationalization of production''; and this, in conditions where other forms of rationalization were limited by scarcity of capital and technical resources, meant above all to raise productivity through a greater physical intensity of labour (see pp. 112-113 above). The campaign to increase the efficiency of the worker was waged on all fronts. Drunkenness, absenteeism and malingering were said to be characteristic of peasants drafted into industry from the countryside rather than of true proletarians; but they seem to have been too prevalent to be curbed by threats of instant dismissal. Side by side with the factory schools which combined vocational training with party indoctrination, the Central Institute of Labour established schools in which young workers received intensive instruction in modern factory techniques. Critics condemned these, as Lenin had once criticized ``Taylorism'' (see p. 25 above), for treating the worker as ``an adjunct of the machine, not a creator of production''. Other forms of incentive were not neglected. What was called ``socialist emulation'' between factories or groups of workers was encouraged by propaganda and by the offer of prizes. The title ``Hero of Labour'', carrying with it certain privileges, was conferred on specially meritorious workers, and the Order of the Red Banner of Labour was created for award to factories, industrial enterprises or workers' collectives. To celebrate the tenth anniversary of the revolution, ``Communist Saturdays'' in imitation of the Communist Saturdays instituted by Lenin in 1919---over-time without pay---were worked in factories and mines in several parts of the USSR. 
% 注：去除了 "(was subjected", "/still unstated", "/should deliver" 等单侧符号。

A device tried at this time was symptomatic of the intense pressure now applied to the workers. On the eve of the tenth anniversary of the revolution in November 1927 the authorities announced a planned transition to a seven-hour working day. This project, hailed as a great achievement of the revolution, was denounced by the opposition as a demagogic attempt to lull the workers into quiescence with the unsubstantial vision of a remote future. But it soon transpired that another end was in view. In order to secure maximum utilization of plant and machinery some factories already worked two shifts in 24 hours. It was now proposed that three continuous shifts of seven hours each should be worked, leaving only a minimum of three hours for cleaning and maintenance. The three-shift system was disliked both by managers and workers, and was at first introduced only in textile factories, which employed almost exclusively women, the lowest-paid category. Incidentally, this involved a total abandonment of the prohibition enacted in the first idealistic days of the revolution, but long honoured in the breach rather than in the observance, on night work for women. During the next two years pronouncements were made from time to time in favour of the extension of the three-shift system throughout industry. But resistance to the tensions and pressures involved was strong. It was noted that, where the system had been introduced, the productivity of the labour employed progressively declined through the working period; and it never appears to have been widely adopted outside the textile industry. 
% 注：去除了极端的乱码 "U4i __ ]employed" 并恢复正常的拼写 "employed"。

Wages, however, remained the focus of the relation of the industrial worker to the employer and to the state. Wages under NEP were in principle fixed by agreement between worker and employer---normally by a collective contract between the trade union and the enterprise or institution concerned. The principle was not affected by the recognition on both sides of a link between productivity and earnings, which was written into the contracts. What radically changed the situation from 1926 onwards was acceptance of the over-riding importance of planning. The wages bill was too vital an element in the economy to be excluded from the calculations of the planners, or to be subject to fluctuations determined by extraneous causes. After much controversy between the trade unions and Vesenkha in which both sides purported to recognize the compatibility of planning with collective contracts, wage-fixing became subject in practice to two distinct processes. First of all, the highest authority---generally the Politburo itself---fixed the total wages fund for the coming year (in an inflationary period some increase in monetary terms was inevitable), and determined what increases should be conceded to which industries. It thus not only fixed the limit of wage payments for which provision had to be made in the plan, but decided which industries were to be encouraged to expand. The second stage was the conclusion of collective contracts, which might be between the trade union central committee and an industry as a whole, or between local trade union committees and individual enterprises. But, since wages had to be kept within limits already laid down, little freedom of negotiation remained; and the discussions of collective contracts were more likely to turn on the conditions of labour or on the ``norms'' of production. 
% 注：去除了 "' ing", "/ an element", "/ of the", "I I to recognize", "I by extraneous", "I contracts", "I Politburo", "I year", "I terms", "I be" 等等大量的侧边划线留下的乱码 "I" 或 "/"。

The original limitation on piece-rates had long ago disappeared; and, where piece-rates were inapplicable, bonuses geared to productivity were a regular part of wage payments. These procedures, which were an essential part of the campaign to gear wages to productivity, required the constant fixing of ``norms'', or the rate for the job. When a general wages increase was granted in the autumn of 1926, Vesenkha began an agitation for a revision of norms. This could be justified in part by measures of rationalization and mechanization, which increased productivity without imposing additional strain on the worker. But to increase norms was more often simply a way to reduce wages. Controversy between Vesenkha and the trade unions raged through 1927, but ended in recognition of the need for a general re-examination of norms. Wage statistics for the period are dispersed, complex and sometimes misleading. In inflationary conditions, unchanged or even increased monetary payments disguised a decline in real wages. It is certain that, whereas the real wages of the worker rose slowly, but steadily, between 1923 and 1927, for several years after 1928 real wages fell, and the worker, no less than other sectors of society, was subjected to the harsh pressures of industrialization, and constricted by the iron hand of the planned economy. 
% 注：去除了 "-y to expand", "^'rose", "t heights" 等乱码。

The role of the trade unions had been a matter of controversy in the early years of the regime. The compromise of NEP had rejected ``the militarization of labour'', and kept the unions formally independent of the state. This independence proved, however, illusory. Under NEP the ``commanding heights'' of industry were firmly held by the state; and it was unthinkable that the trade unions, still ostensibly non-party, but now entirely controlled by the Bolsheviks, should place themselves in opposition to the interests and policies of the workers' state. The first erosion of the independence of the unions came from their commitment to increased productivity. This obliged them to accept responsibility for maintaining labour discipline, and for preventing such ``anarchistic methods'' as strikes and walk-outs. A strike was treated as proof of the failure of the unions to exercise due vigilance and to attend to the needs of the workers. But they could no longer come out unconditionally as supporters of the short-term demands of the workers; their role was rather to serve as mediators, in high-level debates within the party, between these demands and the long-term needs of state industry. At factory level control was in the hands of a ``triangle'' consisting of representatives of the unions, of the management and of the party. But, where the two last were in agreement, the union representative was in a weak position; and the unions were sometimes accused of succumbing to a ``managerial deviation''. 
% 注：去除了 "iparty", "Iplace", "las mediators", "'unions" 等符号。

Moreover, with the rapid expansion of the labour force and of trade union membership, the character itself of the unions underwent a subtle change. The assumption that most factory workers were class-conscious proletarians, with a fair sprinkling of active party workers, was rapidly ceasing to be true. Many of the politically active workers had been promoted to managerial or official positions. Many of the new recruits to industry were peasants fresh from the countryside, who had everything to learn in the way of party doctrine and trade union practices. Much emphasis was laid at this time on the educational role of the unions. But another consequence was a rapid increase in the authority of the leaders, represented by the trade union central council, over the rank-and-file members. 
% 注：去除了 ".new recruits", "[side", ".consequence", "[leaders" 等边缘标点。

The NEP compromise was maintained with growing difficulty from 1922 to 1928---the period of Tomsky's uncontested leadership of the trade unions. The period was one of economic recovery, some of the benefits of which, with the help of the unions, accrued to the workers. It was the advent of planning which inexorably brought about the full integration of the unions into the state machine. The organization and remuneration of labour were an important element in any economic plan. Narkomtrud had by this time become an auxiliary of the trade unions. It was the trade union central council, rather than Narkomtrud, which took its place in major discussions of policy side by side with the economic organs responsible for other elements of the plan. But all alike were subject to the supreme authority of the Politburo, and executed its decisions; by the middle nineteen-twenties high trade union officials were almost invariably party members directly subject to party discipline. As time went on, however, Tomsky and many of his colleagues became increasingly restive under pressures imposed on industrial workers by the plan, and the neglect of time-honoured trade union traditions. It was not altogether paradoxical that the trade unions should have come out in opposition to current policies of industrial expansion. When the party central committee met in July 1928, Tomsky joined Bukharin and Rykov to make up the minority of three in the Politburo who sought to slacken the pace of industrialization. 
% 注：去除了 "[of economic", ". and executed", "f|/high", "II directly", "f..was", "[//have", "l expansion", "'minority" 等大量边缘乱码。

Rapidly increasing investment in industry, and primarily in heavy industry, inflated the demand for both agricultural and industrial products in short supply. The scissors crisis of 1923 had demonstrated the impracticability of leaving the conditions of exchange to be regulated by the free play of the market. That lesson was learned, and control of prices became a permanent item of policy. Control of agricultural prices was in theory exercised through the official collections of agricultural products at fixed prices. But from the winter of 1927-1928 inadequate supplies of grain at official prices were obtained from the producers, largely by methods of compulsion, and were supplemented by purchases at higher prices on the private market. Control of industrial prices was more effective, but presented problems of great complexity. From 1926 onwards, amid the growing pressures of industrialization, price policy was a matter of constant controversy. All major industry now being in state hands, the control of the wholesale prices of industrial products was straightforward enough. Ever since the scissors crisis the policy of holding down industrial prices in order to maintain the link with the peasantry had been firmly embedded in party doctrine. Pleas to raise wholesale prices and thus increase the profitability of state industry came from the opposition in 1926 and 1927, and were indignantly rejected as evidence of the opposition's neglect of the peasant. Successive variants of the five-year plan were based on estimates of reduced prices for industrial goods. 
% 注：去除了 ",A11 major" 中的逗号。

Control of wholesale prices was, however, not matched by an equally effective control of retail prices; and it was frequently pointed out that a regime of strictly controlled wholesale prices combined with floating retail prices merely widened what came to be called ``the wholesale-retail price scissors'', and swelled the undesirable profits of middlemen. Retail prices had been fixed since 1924 for an increasing number of standard commodities (see p. 58 above). These prices, rather unwillingly accepted by state and cooperative shops and selling organizations, were supposed to be obligatory for private traders. Enforcement was difficult and uneven. But police measures, and vigorous propaganda against the unpopular Nepmen, were more successful in the cities than in the countryside in restricting private trade. Prices were fixed without regard to the availability or scarcity of supplies. A decree of July 1926 called for ``a reduction of retail prices of products of state industry in short supply''. During 1927 a series of orders and decrees were issued prescribing a reduction of 10 per cent in the retail prices of standard commodities prevailing on January 1; and, though this ambitious aim was not achieved, some prices were brought down, and many Nepmen driven out of business, in the course of the year. The results of these reductions proved in practice almost wholly illusory. They did not improve the real ability of the peasant or the industrial worker to purchase industrial goods, since they were accompanied by a now chronic shortage of supplies both in town and in country. Price levels were no longer a significant indicator of the economic situation. 
% 注：去除了 "fl in", "I fixed" 等乱码。将 "January i" 修正为 "January 1"。

The year 1927 saw the beginning of a prolonged and progressive decline in standards of living, due to the pressures of industrialization and the absorption of available resources for the planned development of heavy industry. Though the disparity between official prices and black market prices was less extreme for industrial than for agricultural products, this brought little advantage to the consumer, since the shortage of consumer goods was quite as severe as the shortage of foodstuffs. The consumer of whatever category was called to bear a heavy share of the burden of industrialization. The headlong advance of industry, concentrated primarily on the production, not of consumer goods, but of means of production, imposed ever-increasing strains on the peasant, on the worker, on every aspect of the economy. Apathy and resignation rather than active resistance seems to have been the mood of those who carried the heaviest burdens. But the industrializers continued passionately to believe that the achievement was worth the cost, and that the cost would be willingly borne, or could be enforced on the unwilling. 
% 注：去除了 ""The year", "• j .brought", "lof consumer", "J foodstuffs", "4 yt Ntff^ A 6uy otcyx.", "fO. and resignation", "' I ({been" 等等非常严重的乱码区。

During 1928 doubts began to penetrate the Politburo itself. The clash of opinions which occurred at the session of the party central committee in July 1928 (see pp. 125-126 above) turned ostensibly on agricultural policy and the pressure on the peasant. But the underlying issue was the pace of industrialization, by which this policy was dictated. What was significant was the open split which now declared itself in the committee between a majority in the Politburo committed to forced industrialization and a dissentient minority consisting of Rykov, Bukharin and Tomsky, who sought to relax pressures all round by slowing down the pace. At the end of September Bukharin expounded his views in a major article in Pravda entitled ``Notes of an Economist''. Starting from the grain crisis, he launched a full-scale attack on current plans of industrialization, which destroyed the equilibrium between agriculture and industry, and the link with the peasantry established by NEP. Investment in industry was being absurdly and incongruously accelerated in face of material shortages not only of grain, but of industrial products of every kind. Agriculture must be allowed to catch up, and industry developed ``on a basis provided by a rapidly growing agriculture''. Bukharin purported to accept the extent of industrialization already achieved. But the strains were now intolerable, and the pace must not be further increased. He ended by criticizing the ``mad pressure'' contemplated in current drafts of the five-year plan. 
% 注：去除了 ". turned", "///the", "|i( lization", "w'as the", "[and incongruously", "f///Agriculture", "ijll loped" 等标点。

``Notes of an Economist'', the last public pronouncement of opposition to the headlong course of industrialization, and the last rearguard action in defence of NEP, was fiercely attacked both by official economists and by Trotsky and his supporters. Priority for agriculture was no longer an acceptable theme. Bukharin, who at this point departed on holiday for the Caucasus, returned in time for a crucial session of the party central committee in November. Vesenkha under the direction of Kuibyshev continued to demand ever-increasing capital investment in industry. The call ``to catch up with and surpass'' the west had already been made by the industrializers. Stalin took up the theme in a major speech to the committee. Technology, he argued, was ``simply rushing ahead'' in the advanced capitalist countries: ``either we achieve this, or they will destroy us''. He cited Peter the Great, whose feverish construction of factories to meet the needs of defence had been ``an attempt to leap out of the framework of backwardness''. It was the backwardness of the Soviet economy, especially of its agricultural sector, and the isolation of the USSR, which made this ``a matter of life and death for our development''. The committee approved a figure of 1650 million rubles for investment in industry during the year. Bukharin resisted weakly, tendered and then withdrew his resignation, did not challenge a vote, and finally participated in the drafting of the resolution. His defeat was masked by a show of agreement and reconciliation, but was none the less unmistakable. The victory of industrialization was sealed by the completion of the first five-year plan and its submission to the Union Congress of Soviets in May 1929.
% 注：去除了 "/(attacked", "{(his" 等符号。